\section*{Розділ II. Завдання ТПК}
\addcontentsline{toc}{section}{Розділ II. Завдання ТПК}

\subsection*{2.1. Основні завдання}
\addcontentsline{toc}{subsection}{2.1. Основні завдання}
    2.1.1. Основними завданнями ТПК є:

        \begin{enumerate}[label=\alph*)]
            \item Представництво та захист прав і законних інтересів студентів Університету перед адміністрацією Університету та її структурними підрозділами.
            \item Погодження рішень адміністрації Університету з питань, передбачених п. 7.2.5 Положення про ОСС, у порядку, визначеному цим Положенням.
            \item Координація діяльності СР ТПК факультетів/інститутів та СР ТПК гуртожитків.
            \item Забезпечення інформування студентів з питань, що стосуються їхніх прав, обов'язків та інтересів.
            \item Сприяння організації та проведенню виборів до постійних органів студентського самоврядування.
            \item Забезпечення передачі справ, документів та поточних проєктів новообраним постійним органам студентського самоврядування.
        \end{enumerate}

\subsection*{2.2. Фінансове забезпечення}
\addcontentsline{toc}{subsection}{2.2. Фінансове забезпечення}
    2.2.1. ТПК здійснює поточне фінансове управління в межах кошторису, затвердженого КСУ, відповідно до Розділу V Положення про ОСС.

    2.2.2. Голова ТПК забезпечує цільове використання коштів та звітує перед КСУ про фінансову діяльність ТПК.

\subsection*{2.3. Представництво в колегіальних органах}
\addcontentsline{toc}{subsection}{2.3. Представництво в колегіальних органах}
    2.3.1. ТПК забезпечує представництво студентів у колегіальних органах управління Університету на перехідний період відповідно до Положення про ОСС та Закону України ``Про вищу освіту''.
